%% kapa: exam=1 round=32 year=2021 period=1 form=A subject=civil_law subject_ko=민법 seq=25 number=25
%% source: Q-Net 감정평가사 기출문제 (cst00309) artlSeq=5208397 file=extracted/1008/5208397/A형/2021년도 제32회 감정평가사 1차 1교시 A형.pdf
%% conversion: kordoc 4.24.51 -> tools/gen1.py (auto); review=proofread-sonnet; flags=none
\begin{question}{25}
선의취득에 관한 설명으로 옳은 것은? (다툼이 있으면 판례에 따름)
\begin{choices}
  \item 선의취득에 관한 민법 제249조는 저당권의 취득에도 적용된다.
  \item 동산의 선의취득에 필요한 점유의 취득은 현실의 인도뿐만 아니라 점유개정에 의해서도 가능하다.
  \item 선의취득의 요건인 선의·무과실의 판단은 동산의 인도 여부와 관계없이 물권적 합의가 이루어진 때를 기준으로 한다.
  \item 도품·유실물에 관한 민법 제251조는 선의취득자에게 그가 지급한 대가의 변상시까지 취득물의 반환청구를 거부할 수 있는 항변권만을 인정한다는 취지이다.
  \item 제3자에 대한 목적물반환청구권을 양수인에게 양도하고 지명채권 양도의 대항요건을 갖추면 동산의 선의취득에 필요한 점유의 취득요건을 충족한다.
\end{choices}
\end{question}
\answer{5}
